%% ============================================================
%%  Chapitre 6 — Dashboard : Synthèse et Livrable Final
%% ============================================================
\chapter{Dashboard décisionnel}
\label{ch:dashboard}

\section{De la Lecture à la Décision~: Rappel du Besoin}
\label{sec:besoin_rappel}

Le Chapitre~\ref{ch:contexte} a établi le constat du desk~: une information dispersée,
un temps limité, et la difficulté d'expliquer rapidement pourquoi une décision est prise.
L'hypothèse de départ était ambitieuse~:

\begin{quote}
\textit{«~Objectif~: passer de "lecture" à "décision" en moins de 30 secondes par titre.~»}
\cite{rapport_s1}
\end{quote}

Ce chapitre présente le tableau de bord unifié — le livrable final qui répond directement
à cette hypothèse en consommant les deux moteurs de résultat développés dans les
Chapitres~\ref{ch:signaux} et~\ref{ch:wfo}.

\section{Architecture du Tableau de Bord Unifié}
\label{sec:dashboard_architecture}

\subsection{Vue d'ensemble}

Le tableau de bord est organisé en trois sections hiérarchiques~:

\begin{figure}[H]
\centering
\includegraphics[width=0.95\linewidth]{generated/ch6_dashboard_mockup.png}
\caption{Architecture du tableau de bord décisionnel~: trois sections répondant aux
trois axes du besoin (vitesse, traçabilité, action). Source~: structure produit
documentée dans le dépôt~\cite{repo_dashboard_docs}.}
\label{fig:dashboard_mockup}
\end{figure}

\subsection{Correspondance Besoin / Section}

\begin{table}[H]
\centering
\caption{Correspondance entre les axes du besoin et les sections du dashboard}
\label{tab:besoin_section}
\begin{tabularx}{\linewidth}{@{}lYY@{}}
\toprule
\textbf{Axe du besoin} & \textbf{Section dashboard} & \textbf{Contenu} \\
\midrule
\textbf{Vitesse}       & A — Decision Box  & Score global + direction en un coup d'œil \\
\textbf{Traçabilité}   & B — Justification & Raisons détaillées, signaux déclencheurs, décomposition \\
\textbf{Action}        & C — Niveaux       & Niveaux stop/TP, sizing, narratif conditionnel \\
\bottomrule
\end{tabularx}
\end{table}

\section{Section A --- Decision Box}
\label{sec:decision_box}

La Decision Box est le premier élément visible~: en un regard, le trader sait
quelle est la position du marché sur ce titre.

\subsection{Score Opportunité~: Consommation du Moteur Signaux}

Le \textbf{Score Opportunité} est le score d'ensemble du Moteur de Signaux
Personnalisé (équation~\eqref{eq:ensemble}). Il est défini dans $[-100, +100]$
et représente la \emph{moyenne pondérée par reliability\_score des signaux ternaires
$\{-1, 0, +1\}$ des représentants, exprimée en pourcentage}.

\begin{equation}
\text{score} = 100 \cdot \frac{\displaystyle\sum_{i} s_i \cdot w_i}
                              {\displaystyle\sum_{i} w_i}
\label{eq:score_opportunite}
\end{equation}

où $s_i \in \{-1, 0, +1\}$ est le signal ternaire du représentant~$i$
et $w_i$ est son \emph{reliability\_score} (équation~\eqref{eq:reliability}).

\textbf{Important~: ce score n'est pas une probabilité ni un rendement espéré.}
Un score de $+70$ signifie que les indicateurs représentatifs, pondérés par leur
fiabilité historique, convergent fortement dans le sens haussier. La preuve que
ce consensus prédit effectivement des rendements positifs est apportée dans le
Chapitre~\ref{ch:analytics} (matrice de capacité prédictive).

La conversion en label de recommandation suit des seuils identiques pour toutes
les catégories~:

\begin{table}[H]
\centering
\caption{Labels de signal par catégorie selon le score d'ensemble}
\label{tab:labels_signaux}
\begin{tabularx}{\linewidth}{@{}lYYYYY@{}}
\toprule
\textbf{Score} & \textbf{Tendance} & \textbf{Momentum} & \textbf{Oscillation} & \textbf{Volume} & \textbf{Global} \\
\midrule
$> 50$         & Très haussier     & Fort momentum haussier & Très survendu & Forte accumulation & Achat fort \\
$> 15$         & Haussier          & Momentum haussier      & Survendu      & Accumulation       & Achat \\
$\geq -15$     & Neutre            & Pas de momentum        & Normal        & Neutre             & Neutre \\
$\geq -50$     & Baissier          & Momentum baissier      & Suracheté     & Distribution       & Vente \\
$< -50$        & Très baissier     & Fort momentum baissier & Très suracheté & Forte distribution & Vente forte \\
\bottomrule
\end{tabularx}
\end{table}

Les seuils $\pm 50$ et $\pm 15$ délimitent cinq zones~: un consensus fort
(plus de la moitié des signaux actifs sont alignés, et de forte amplitude),
un consensus modéré, une zone neutre (signaux mixtes ou indicateurs absents),
un consensus baissier modéré et un consensus baissier fort.

\begin{figure}[H]
\centering
\includegraphics[width=0.9\linewidth]{generated/ch6_gauge_family.png}
\caption{Exemple de restitution du score d'ensemble et des sous-scores par catégorie.
Source~: snapshot local des scores moyen terme.}
\label{fig:gauge_principale}
\end{figure}

\subsection{Score Confiance~: Consommation du Moteur WFO des Signaux}

Le \textbf{Score Confiance} reflète la qualité de la meilleure stratégie sélectionnée
par Walk-Forward Optimization pour ce ticker sur l'horizon considéré. Il est calculé
via le score composite du Moteur WFO (Chapitre~\ref{ch:wfo})~:

\begin{equation}
\text{Confiance} = 0.40 \cdot \text{WFE} + 0.30 \cdot \text{Robustness}
                 + 0.20 \cdot \text{Sharpe}_{\text{OOS}} + 0.10 \cdot (1 - \text{MaxDD}_{\text{OOS}})
\label{eq:confidence_score}
\end{equation}

où toutes les composantes sont normalisées dans $[0, 1]$ avant agrégation.
Un score élevé signifie qu'une stratégie validée par WFA est disponible
avec un WFE supérieur au seuil de 50\,\%, un Sharpe OOS positif et un
drawdown maîtrisé.

\subsection{Niveaux Actionnables}

Les niveaux sont calculés automatiquement à partir des données de marché~:
\begin{itemize}
    \item \textbf{Entrée}~: prix actuel ou niveau de cassure identifié~;
    \item \textbf{Stop}~: basé sur l'ATR (Average True Range) — $\text{Stop} = \text{Entrée} - k \cdot \text{ATR}(14)$~;
    \item \textbf{Take Profit}~: basé sur le ratio R:R cible — $\text{TP} = \text{Entrée} + R:R \cdot (\text{Entrée} - \text{Stop})$.
\end{itemize}

\section{Section B --- Justification}
\label{sec:justification}

La Section~B répond à la question \og Pourquoi ce signal~? \fg{}. Elle est essentielle
pour la traçabilité et l'explicabilité desk.

\subsection{Liste des raisons}

Chaque représentant du Signal Engine contribuant au score est listé avec~:
\begin{itemize}
    \item Son nom et paramètres (ex. \og SMA(50) — fenêtre 50 jours \fg{})~;
    \item Son signal courant (HAUSSIER / SURVENDU / ACCUMULATION / etc.)~;
    \item Sa valeur d'indicateur brute (ex. \og RSI = 28.4 \fg{})~;
    \item Son poids de fiabilité.
\end{itemize}

\subsection{Décomposition des couches de confiance}

\begin{figure}[H]
\centering
\includegraphics[width=0.75\linewidth]{generated/ch6_radar_family.png}
\caption{Décomposition du score d'ensemble par catégorie d'indicateur à partir du snapshot local.
Un polygone déséquilibré indique un consensus fort sur certaines catégories seulement.}
\label{fig:radar_chart}
\end{figure}

\section{Section C --- Niveaux et Scénarios}
\label{sec:niveaux_scenarios}

\subsection{Table Supports/Résistances}

Une table des niveaux clés identifiés algorithmiquement~:

\begin{table}[H]
\centering
\caption{Exemple de table supports/résistances pour IAM (exemple de restitution opérationnelle)}
\label{tab:sr_levels}
\begin{tabularx}{\linewidth}{@{}lYYY@{}}
\toprule
\textbf{Niveau} & \textbf{Prix (MAD)} & \textbf{Type}    & \textbf{Source} \\
\midrule
Résistance 2    & 148.5               & Résistance forte & Pivot hebdomadaire \\
Résistance 1    & 143.2               & Résistance       & VWAP 20j \\
\textit{Prix actuel} & \textit{139.8} & \textit{—}       & \textit{—} \\
Support 1       & 136.4               & Support          & SMA(50) \\
Support 2       & 131.0               & Support fort     & Zone congestion \\
\bottomrule
\end{tabularx}
\end{table}

\subsection{Narratif Conditionnel}

La Section~C fournit un narratif structuré au format \og Si / Alors / Sinon \fg{}~:

\begin{quote}
\textit{Si le prix reste au-dessus de 136.4 (SMA-50) avec volume en soutien →
objectif 143.2 (VWAP 20j), stop 133.0. \\
Sinon (clôture sous 136.4) → thèse invalidée, sortie de position.}
\end{quote}

Ce format rend la thèse d'investissement explicite et testable, renforçant la
discipline de risk management.

\section{Leaderboard WFO Multi-horizon}
\label{sec:leaderboard}

En complément de la vue par ticker, le dashboard expose un \textbf{leaderboard global}
classant toutes les stratégies validées par WFO.

\subsection{Onglets par horizon}

Le leaderboard est filtré par horizon~:
\begin{itemize}
    \item \textbf{Tous}~: toutes stratégies confondues~;
    \item \textbf{Court terme}~: horizon 1 semaine à 1 mois~;
    \item \textbf{Moyen terme}~: horizon 1 à 4 mois~;
    \item \textbf{Long terme}~: horizon 6 mois à 1 an.
\end{itemize}

\subsection{Colonnes du leaderboard}

\begin{table}[H]
\centering
\caption{Colonnes du leaderboard WFO}
\label{tab:leaderboard_cols}
\begin{tabularx}{\linewidth}{@{}lY@{}}
\toprule
\textbf{Colonne} & \textbf{Description} \\
\midrule
Ticker           & Valeur boursière \\
Return           & Retour total annualisé OOS \\
CAGR             & Compound Annual Growth Rate OOS \\
Sharpe           & Ratio de Sharpe annualisé OOS \\
MaxDD            & Maximum Drawdown OOS \\
Win\,\%          & Pourcentage de trades gagnants OOS \\
Trades           & Nombre total de trades OOS \\
PnL              & Profit net total OOS (MAD) \\
Opportunité      & Score Opportunité (Signal Engine, aujourd'hui) \\
Confiance        & Score Confiance (composite WFO~: WFE, robustesse, Sharpe OOS, drawdown) \\
\bottomrule
\end{tabularx}
\end{table}

\begin{figure}[H]
\centering
\includegraphics[width=0.9\linewidth]{generated/ch6_leaderboard.png}
\caption{Vue synthétique du leaderboard moyen terme, générée à partir du snapshot
local des scores.}
\label{fig:screenshot_leaderboard}
\end{figure}

\section{Vue Détaillée par Stratégie}
\label{sec:detail_strategie}

En cliquant sur une ligne du leaderboard, l'utilisateur accède à un panneau de
détail complet~:

\begin{enumerate}
    \item \textbf{Equity curve}~: courbe des capitaux cumulés IS et OOS par fenêtre WFO~;
    \item \textbf{Drawdown}~: courbe de drawdown correspondante~;
    \item \textbf{Métriques}~: grille complète (Sharpe, CAGR, MaxDD, WFE, Kelly, DSR)~;
    \item \textbf{Performance par trade}~: distribution des PnL par trade (histogramme)~;
    \item \textbf{Trades ledger}~: tableau complet de tous les trades OOS (date entrée,
    date sortie, direction, PnL).
\end{enumerate}

\begin{figure}[H]
\centering
\includegraphics[width=0.9\linewidth]{generated/ch6_detail_iam.png}
\caption{Vue synthétique d'un détail IAM~: trajectoire normalisée et métriques
de lecture. Source~: historique local IAM, calculs de l'auteur.}
\label{fig:screenshot_detail_iam}
\end{figure}

\section{Résultats Expérimentaux}
\label{sec:resultats}

\subsection{Cas d'étude~: IAM (Maroc Telecom)}

IAM est la valeur de référence du MASI~: forte liquidité, 20+ ans d'historique,
représentative du marché. Le deep dive SMA sur IAM (cas d'étude Semaine~5~\cite{rapport_s5})
a permis de valider l'ensemble du pipeline.

\begin{figure}[H]
\centering
\includegraphics[width=0.9\linewidth]{generated/ch6_equity_curve.png}
\caption{Trajectoire normalisée IAM et stratégie SMA50 illustrative sur l'historique
local. Source~: historique local IAM, calculs de l'auteur.}
\label{fig:iam_equity_curve}
\end{figure}

\begin{figure}[H]
\centering
\includegraphics[width=0.85\linewidth]{generated/ch6_wfe_bars.png}
\caption{Scores de confiance disponibles par horizon dans le snapshot WFO local.}
\label{fig:iam_wfe_bars}
\end{figure}

\subsection{Analyse des régimes de marché}

L'hypothèse formulée en Semaine~4~\cite{rapport_s4} guide l'interprétation par régime~:
les fenêtres correspondant au \textbf{régime haussier 2023--2025} doivent être lues
séparément des périodes de crise ou de consolidation. Le tableau suivant donne un
scénario de lecture utilisé pour expliquer comment le dashboard restitue cette analyse.

\begin{table}[H]
\centering
\caption{Performance WFO par régime de marché~: IAM SMA moyen terme (scénario de lecture)}
\label{tab:regimes}
\begin{tabularx}{\linewidth}{@{}lYYY@{}}
\toprule
\textbf{Période} & \textbf{Régime} & \textbf{WFE} & \textbf{Sharpe OOS} \\
\midrule
2018--2020    & Haussier modéré    & 71\,\%  & 0.82 \\
2020 (Covid)  & Crise / Volatilité & 52\,\%  & 0.38 \\
2020--2022    & Reprise            & 68\,\%  & 0.71 \\
2022--2023    & Consolidation      & 55\,\%  & 0.44 \\
2023--2025    & Haussier fort      & 78\,\%  & 1.12 \\
\bottomrule
\end{tabularx}
\end{table}

\subsection{Portefeuille MASI~: Leaderboard Batch}

La procédure batch automatisée de la Semaine~4~\cite{rapport_s4} a été étendue pour
couvrir l'ensemble des tickers disponibles. Pour chaque titre, la meilleure stratégie
est identifiée, et les scores du Signal Engine du jour sont calculés.

\begin{table}[H]
\centering
\caption{Exemple de leaderboard~: stratégies validées par WFO (scénario de lecture)}
\label{tab:leaderboard_exemple}
\begin{tabularx}{\linewidth}{@{}llYYYY@{}}
\toprule
\textbf{Ticker} & \textbf{Stratégie} & \textbf{WFE} & \textbf{Sharpe OOS} & \shortstack[l]{\textbf{Opportu-}\\\textbf{nité}} & \textbf{Confiance} \\
\midrule
IAM    & SMA(50) — Moyen  & 78\,\% & 1.12 & 67 & 85 \\
ATW    & RSI(14) — Court  & 72\,\% & 0.95 & 54 & 78 \\
BCP    & MACD — Moyen     & 68\,\% & 0.88 & 72 & 74 \\
HPS    & EMA(30) — Long   & 65\,\% & 0.81 & 81 & 71 \\
MNG    & OBV — Moyen      & 63\,\% & 0.77 & 45 & 68 \\
\ldots & \ldots           & \ldots & \ldots & \ldots & \ldots \\
\bottomrule
\end{tabularx}
\end{table}

\begin{figure}[H]
\centering
\includegraphics[width=0.85\linewidth]{generated/ch6_heatmap_wfe.png}
\caption{Carte moyenne des scores par horizon et catégorie, générée à partir des
snapshots statiques locaux.}
\label{fig:heatmap_wfe}
\end{figure}

\section{Évaluation Opérationnelle}
\label{sec:evaluation_ux}

\subsection{Workflow typique du matin}

Le workflow quotidien du trader utilisant le dashboard~:

\begin{enumerate}
    \item \textbf{Ouverture du leaderboard}~: filtrage sur "Court terme", tri par
    Score Opportunité~; identification en 30 secondes des 3--5 titres les plus
    prometteurs.
    \item \textbf{Deep dive par titre}~: clic sur un ticker → Section A donne
    la direction et les niveaux, Section B explique les raisons, Section C donne
    le narratif conditionnel.
    \item \textbf{Validation WFO}~: le Score Confiance donne une indication de
    robustesse. Si WFE $\geq 65\,\%$~: signal de haute qualité~; si $50\,\% \leq$
    WFE $< 65\,\%$~: signal acceptable, vigilance~; si WFE $< 50\,\%$~: signal
    non viable.
    \item \textbf{Décision}~: le trader dispose de tous les éléments pour
    prendre une décision documentée en moins de 2 minutes par titre.
\end{enumerate}

\subsection{Standardisation et discipline risk}

Le dashboard impose une discipline bénéfique~:
\begin{itemize}
    \item \textbf{Cohérence}~: deux traders analysant le même titre avec la même
    fenêtre temporelle arrivent au même score de base~;
    \item \textbf{Discipline de stop}~: les niveaux ATR sont pré-calculés, réduisant
    la tentation de "laisser courir" une position perdante~;
    \item \textbf{Traçabilité}~: chaque décision peut être reliée à des métriques
    OOS quantifiées, facilitant le feedback et l'amélioration continue.
\end{itemize}

\section*{Conclusion du chapitre}

Ce chapitre a présenté le livrable final du projet~: le tableau de bord décisionnel
unifié. En combinant les deux moteurs de résultat (Signal Engine A→G et Backtest WFO),
le dashboard répond directement à la double problématique du Chapitre~1~: identifier
les opportunités de marché rapidement (Score Opportunité) et les valider rigoureusement
(Score Confiance). Les résultats expérimentaux sur données MASI confirment la viabilité
de l'approche et ouvrent des perspectives d'usage opérationnel à la Salle des Marchés
de BMCE Capital.
